\chapter{Polímeros Termoestables}
\label{cap:termoestables}

Los polímeros termoestables se distinguen estructuralmente de los termoplásticos por poseer una red tridimensional insoluble e infusible unida covalentemente mediante reacciones de curado tridimensionales. Estos materiales ofrecen una excelente estabilidad dimensional, gran resistencia mecánica a la deformación bajo cargas térmicas altas y nula fluencia (\textit{creep}).

\section{Química del Curado y Formación de Redes}

\subsection{Resinas Epoxi}
Las resinas epoxídicas son el estándar de oro en adhesivos y matrices estructurales aeroespaciales. El curado se basa en la reacción de apertura del anillo oxirano reactivo terminal empleando agentes nucleofílicos multifuncionales como las diaminas primarias. Los átomos de nitrógeno reaccionan con cuatro anillos epoxídicos individuales, actuando como puntos de unión de red molecular.

\section{Termodinámica del Curado: Gelificación}

Durante el curado térmico de una mezcla reactiva líquida compuesta por monómeros multifuncionales, ocurre una transición termodinámica llamada gelificación o punto de gel. 

\subsection{El Punto de Gel estadístico (Flory-Stockmayer)}
Flory y Stockmayer dedujeron la ecuación estadística que predice la conversión de grupos funcionales crítica ($p_c$) necesaria para que una red infinitamente ramificada insoluble aparezca en el reactor:
\begin{equation}
    p_c = \frac{1}{f - 1}
    \label{eq:punto_gel}
\end{equation}
donde $f$ representa la funcionalidad del reactivo homólogo (número de grupos reactivos por molécula). Por ejemplo, para un monómero trifuncional ($f = 3$), el punto de gel ocurre exactamente a una conversión de grupos funcionales del $p_c = 0.50$ (50.0\%). Justo en el punto de gel, la viscosidad del medio de reacción se vuelve infinita y el material pierde por completo su capacidad de fluir.

La vitrificación es otra transición física relevante que ocurre cuando la temperatura de transición vítrea de la red reactiva en crecimiento iguala a la temperatura de curado del reactor ($T_g = T_{\text{curado}}$), congelando la cinética difusiva de reacción.
